\documentclass[a4paper]{article}
% Español (México) con XeLaTeX: fontspec para fuentes Unicode, polyglossia
% para la división silábica y los nombres del documento en la variante mexicana.
\usepackage{fontspec}
\usepackage{polyglossia}
\setdefaultlanguage[variant=mexican]{spanish}
% Los nombres chinos van como «pinyin (original)»: xeCJK da a los caracteres
% Han su propia fuente; sin ella XeLaTeX los omite con un aviso.
\usepackage{xeCJK}
\setCJKmainfont{FandolSong-Regular.otf}
% polyglossia no traduce los nombres de listings.
\AtBeginDocument{\providecommand{\lstlistingname}{}\renewcommand{\lstlistingname}{Listado}%
  \providecommand{\lstlistlistingname}{}\renewcommand{\lstlistlistingname}{Índice de listados}}
\usepackage{amsmath, amssymb}
\usepackage{graphicx}
\usepackage[margin=2.5cm]{geometry}
\usepackage[most]{tcolorbox}
\usepackage{etoolbox}
\usepackage{listings}
\usepackage{hyperref}
\usepackage{booktabs}
\usepackage{float}
\newtcolorbox{knowledgebox}[1]{enhanced, colback=blue!5!white, colframe=blue!75!black, colbacktitle=blue!75!black, coltitle=white, fonttitle=\bfseries, title={#1}, attach boxed title to top left={yshift=-2mm, xshift=2mm}, boxrule=1pt, sharp corners}
\newtcolorbox{importantbox}[1]{enhanced, colback=yellow!10!white, colframe=yellow!80!black, colbacktitle=yellow!80!black, coltitle=black, fonttitle=\bfseries, title={#1}, sharp corners}
\newtcolorbox{warningbox}[1]{enhanced, colback=red!5!white, colframe=red!75!black, colbacktitle=red!75!black, coltitle=white, fonttitle=\bfseries, title={#1}, sharp corners}
\lstset{language=Python, basicstyle=\ttfamily\small, keywordstyle=\color{blue}, stringstyle=\color{red!60!black}, commentstyle=\color{green!60!black}, breaklines=true, frame=single, numbers=left, numberstyle=\tiny\color{gray}, captionpos=b, extendedchars=true}
\newcommand{\notetitle}{DPO: de Reward Model a optimización directa de preferencias}
\newcommand{\noteauthors}{Elaboradas a partir de materiales públicos del curso}
\newcommand{\notedate}{2025}
\newcommand{\videochannel}{Wudaokou Nashi (五道口纳什)}
\newcommand{\videopublishdate}{2025}
\newcommand{\videoduration}{aproximadamente 20 minutos}
\newcommand{\videourl}{}
\newcommand{\videocoverpath}{cover.jpg}
\newcommand{\repourl}{https://github.com/hqhq1025/ai-course-notes}

% Footer with repo info
\usepackage{fancyhdr}
\pagestyle{fancy}
\fancyhf{}
\fancyfoot[C]{\thepage}
\fancyfoot[R]{\footnotesize\color{gray}\href{\repourl}{AI Course Notes}}
\renewcommand{\headrulewidth}{0pt}
\renewcommand{\footrulewidth}{0pt}

\begin{document}
\begin{titlepage}\centering
{\Large Notas del curso\par}\vspace{1.2cm}{\huge\bfseries \notetitle\par}\vspace{0.8cm}{\large \noteauthors\par}\vspace{0.3cm}{\large \notedate\par}\vspace{1.2cm}
\ifdefempty{\videocoverpath}{}{\includegraphics[width=0.82\textwidth,height=0.45\textheight,keepaspectratio]{\videocoverpath}\par}
\vfill
\begin{tcolorbox}[width=0.9\textwidth, colback=black!2!white, colframe=black!60, sharp corners]
\textbf{Autor o canal del video}: \videochannel\par\textbf{Fecha de publicación}: \videopublishdate\par\textbf{Duración del video}: \videoduration\par
\ifdefempty{\videourl}{}{\textbf{Enlace del video}: \href{\videourl}{\nolinkurl{\videourl}}\par}\par
\textbf{Página del proyecto}: \href{\repourl}{\nolinkurl{\repourl}}\par
\end{tcolorbox}
\end{titlepage}
\tableofcontents\newpage

\section{Introducción}
Con la preparación de los dos episodios anteriores sobre el Reward Model y la perspectiva distribucional, este episodio presenta DPO (Direct Preference Optimization).

\begin{importantbox}{Idea central de DPO}
DPO combina las dos etapas de RLHF (primero entrenar el RM, luego usar el RM para hacer RL) en un solo paso de optimización directa, evitando un Reward Model explícito.
\end{importantbox}

\section{De RLHF a DPO}
\subsection{Las dos etapas de RLHF}
\begin{enumerate}
  \item Recopilar datos de preferencia $\to$ entrenar el Reward Model
  \item Usar el RM para construir una pérdida sustituta $\to$ optimizar la política con PPO
\end{enumerate}

\subsection{Derivación de DPO}
El conocimiento clave de DPO: existe una relación de forma cerrada entre la política óptima $\pi^*$ y la reward function. Aprovechando esta relación, se puede optimizar la política directamente a partir de los datos de preferencia, sin un RM explícito.

$$\mathcal{L}_{\text{DPO}} = -\mathbb{E}\left[\log \sigma\left(\beta \log \frac{\pi_\theta(y_w|x)}{\pi_{\text{ref}}(y_w|x)} - \beta \log \frac{\pi_\theta(y_l|x)}{\pi_{\text{ref}}(y_l|x)}\right)\right]$$

\subsection{Resumen de la sección}
DPO comprime, mediante una derivación matemática, las dos etapas de RLHF en un solo paso: una simplificación elegante.

\newpage
\section{Comparación de métodos: simplificar el proceso no significa dejar de tener supuestos}
El atractivo de DPO radica en que sustituye el proceso de dos etapas de RLHF por una cadena de optimización más corta, pero esto no significa que carezca de premisas. Es válido porque aceptamos una relación específica entre la política óptima y el reward, y porque estamos dispuestos a usar la reference policy y la restricción de KL para mantener estable el entrenamiento.

\begin{knowledgebox}{Tres puntos que hay que conservar al entender DPO}
\begin{itemize}
  \item Lo que elimina es el RM explícito, no el modelado de preferencias en sí
  \item La reference policy sigue desempeñando el papel de estabilizador en el objetivo
  \item La calidad de los datos de preferencias sigue determinando de manera directa el límite superior del entrenamiento
\end{itemize}
\end{knowledgebox}

\subsection{Resumen de la sección}
La «elegancia» de DPO proviene de la sencillez de su derivación, pero en la práctica sigue siendo necesario atender la calidad de los datos, la elección de la reference y el ruido en las preferencias; no debe entenderse como un sustituto sin costo.

\newpage
\section{Síntesis y ampliación}
\begin{enumerate}
  \item DPO evita el RM explícito y optimiza la política directamente a partir de datos de preferencias
  \item Es matemáticamente equivalente a RLHF con restricción de KL
  \item Es simple de implementar y su desempeño es comparable o superior al de RLHF
\end{enumerate}
\end{document}

